\section*{الفصل \ref{ch:b2:mantle-melting} --- انصهار الوشاح}

\begin{solution}{exo:b2:mantle-melting:1}
(أ) أكثر من \qty{90}{\percent} من الأوليفين: دونيت. (ب) أوليفين وأورثوبيروكسين
دون كلينوبيروكسين: هارزبورجيت. (ج) أكثر من \qty{40}{\percent} من الأوليفين مع
البيروكسينين كليهما: لرزوليت. (د) أقل من \qty{40}{\percent} من الأوليفين: بيروكسينيت
(ويبستريت أوليفيني، فيه البيروكسينان كلاهما). (ه) $57/98 = \qty{58.2}{\percent}$،
$28/98 = \qty{28.6}{\percent}$، $13/98 = \qty{13.3}{\percent}$: لرزوليت.
\end{solution}

\begin{solution}{exo:b2:mantle-melting:2}
(أ) خط السيولة عند \qty{60}{\percent} من الأنورثيت هو \qty{1390}{\celsius}، في
حقل الأنورثيت: فيتبلور الأنورثيت أولًا. (ب) يظهر الديوبسيد عند
الأوتكتيك، \qty{1270}{\celsius}، حين يحوي السائل \qty{42}{\percent} من الأنورثيت.
(ج) أول سائل لأيّ خليط من المعدنين هو السائل الأوتكتيكي،
\qty{42}{\percent} من الأنورثيت، عند \qty{1270}{\celsius}.
\end{solution}

\begin{solution}{exo:b2:mantle-melting:3}
(أ) $\num{4.0e-5}\times1623/(3300\times1000) = \qty{1.97e-8}{K/Pa} =
\qty{19.7}{K/GPa}$. (ب) $\num{e9}/(3300\times9.81) = \qty{3.09e4}{m}$، أي
\qty{30.9}{km} لكل غيغاباسكال؛ $19.7/30.9 = \qty{0.64}{K/km}$، أي أقل بنحو خمسين
مرة من \qty{30}{K/km} قرب السطح، حيث تنتقل الحرارة بالتوصيل لا
بالحمل.
\end{solution}

\begin{solution}{exo:b2:mantle-melting:4}
(أ) $T_{\mathrm s}(0) = \qty{1085.7}{\celsius}$؛ $T_{\mathrm s}(1) = 1085.7 +
132.9 - 5.1 = \qty{1213.5}{\celsius}$؛ $T_{\mathrm s}(2) = 1085.7 + 265.8 - 20.4 =
\qty{1331.1}{\celsius}$؛ $T_{\mathrm s}(3) = 1085.7 + 398.7 - 45.9 =
\qty{1438.5}{\celsius}$. (ب) $132.9 - 2\times5.1\times2 = \qty{112.5}{K/GPa}$.
(ج) عند \qty{3}{GPa} تكون \qty{1400}{\celsius} دون خط الصلابة،
\qty{1438.5}{\celsius}؛ وحين يصعد الصخر إلى \qty{2}{GPa} على امتداد التدرّج الكظوم لا يبرد إلا
بنحو \qty{20}{\celsius}، إلى قرابة \qty{1380}{\celsius}، بينما ينخفض خط الصلابة إلى
\qty{1331.1}{\celsius}: فيصبح الصخر فوق خط صلابته وينصهر جزء منه.
\end{solution}

\begin{solution}{exo:b2:mantle-melting:5}
(أ) $1350 + (1277.5 - 1350)\times(30 - 20)/(40 - 20) = \qty{1314}{\celsius}$.
(ب) يقع التركيب في جهة الديوبسيد من الأوتكتيك: الديوبسيد.
(ج) يبلغ السائل الأوتكتيك وفيه \qty{42}{\percent} من الأنورثيت؛ ويبقى الأنورثيت
في السائل، فكسر السائل $30/42 = 0.714$، و$1 - 0.714 =
\qty{29}{\percent}$ من النظام تبلور ديوبسيدًا. (د) يتبلور
كل شيء: \qty{70}{\percent} من الديوبسيد و\qty{30}{\percent} من الأنورثيت، أي
التركيب الكلي.
\end{solution}

\begin{solution}{exo:b2:mantle-melting:6}
(أ) الفورستريت: $60.08/(2\times40.30 + 60.08) = 60.08/140.68 =
\qty{42.71}{\percent}$ من \ce{SiO2}؛ الإنستاتيت: $60.08/100.38 = \qty{59.85}{\percent}$.
(ب) $0.975\times59.85 + 2.5 = \qty{60.85}{\percent}$. (ج) يقع الإنستاتيت بين
الفورستريت والسائل: فكسر الفورستريت $(60.85 - 59.85)/(60.85 -
42.71) = 1.00/18.14 = \qty{5.5}{\percent}$، وهي قيمة بوين وأندرسن. (د)
المصهور أغنى بالسيليكا من الأورثوبيروكسين الذي جاء منه،
والبقايا تكسب أوليفينًا.
\end{solution}

\begin{solution}{exo:b2:mantle-melting:7}
(أ) $5.1P^2 - 112.9P + 214.3 = 0$: $P_0 = (112.9 - \sqrt{12746.4 -
4371.7})/10.2 = (112.9 - 91.51)/10.2 = \qty{2.10}{GPa}$، على عمق $2.097\times30.9 =
\qty{64.8}{km}$. (ب) $5.1P^2 - 112.9P + 364.3 = 0$: $P_0 = (112.9 -
72.90)/10.2 = \qty{3.92}{GPa}$، على عمق \qty{121}{km}. (ج) $T_{\mathrm s}'(2.10) = 132.9
- 21.4 = \qty{111.5}{K/GPa}$، ومنه $\Delta P_0 = 10/(111.5 - 20) = \qty{0.109}{GPa}$
لكل \qty{10}{\celsius}، أي \qty{3.4}{km}.
\end{solution}

\begin{solution}{exo:b2:mantle-melting:8}
(أ) $T_{\mathrm s} = \qty{1331.1}{\celsius}$، $T_{\mathrm l} = 1475.0 + 160.0 -
12.8 = \qty{1622.2}{\celsius}$؛ $T = 1331.1 + 0.10\times291.1 =
\qty{1360.2}{\celsius}$، أي \qty{1633}{K}. (ب) $(\partial T/\partial P)_F = 112.5
+ 0.10\times(67.2 - 112.5) = \qty{108.0}{K/GPa}$؛ $(\partial T/\partial F)_P =
T_{\mathrm l} - T_{\mathrm s} = \qty{291.1}{K}$. (ج) $c_p/T = 0.612$؛ ومع $\alpha/\rho
= \qty{12.1}{J/(kg.K)}$ لكل غيغاباسكال، يكون البسط $0.612\times108.0 - 12.1
= 54.0$ والمقام $300 + 0.612\times291.1 = 478.2$: $-\mathrm dF/\mathrm
dP = 0.113$، أي \qty{11.3}{\percent} لكل غيغاباسكال. (د) $54.0/178.2 = 0.303$،
أي \qty{30}{\percent} لكل غيغاباسكال: فالحرارة الكامنة تخفض الإنتاجية بعامل
2.7.
\end{solution}

\begin{solution}{exo:b2:mantle-melting:9}
(أ) $43\times1 = \qty{43}{\celsius}$؛ $43\times5^{0.75} = 43\times3.34 =
\qty{144}{\celsius}$؛ $43\times13^{0.75} = 43\times6.85 = \qty{294}{\celsius}$.
(ب) $1213.5 - 294 = \qty{919}{\celsius}$. (ج) \qty{125}{ppm} تساوي
\qty{0.0125}{\percent}، ومنه $X = 0.0125/0.01 = \qty{1.25}{\percent}$ و$\Delta T =
43\times1.25^{0.75} = 43\times1.18 = \qty{51}{\celsius}$.
\end{solution}

\begin{solution}{exo:b2:mantle-melting:10}
(أ) عرض القاعدة $2H$ ويعبرها الوشاح بالسرعة $U$: أي $2UH$ لكل وحدة طول
من الظهرة ولكل وحدة زمن. ويمتدّ كل ضلع على الارتفاع $H$ ويُعبَر أفقيًا بالسرعة
$U$: أي $UH$ لكل منهما، و$2UH$ في المجموع. ولو كان الضلعان بزاوية $\theta$ لكان عرض القاعدة
$2H/\tan\theta$ والتدفّق الداخل $2UH/\tan\theta$، ولا يساوي الخارج إلا إذا
كان $\tan\theta = 1$: ففي هذا النموذج تحدّد الزاويةَ مصونيةُ الكتلة. (ب) مع
إنتاجية ثابتة يكون \omterm{def:b2:mantle-melting:fraction}{كسر المصهور} عند الرأس $F_{\max} = \Pi P_0 = \Pi\rho
g z_0$، و$H = z_0$: $h_c = F_{\max}z_0/2 = \Pi\rho g z_0^2/2$. (ج) في كل سنة
\qty{50}{mm} من صفيحة جديدة، أي $\qty{0.05}{m}\times\qty{1000}{m}\times\qty{7100}{m} =
\qty{3.6e5}{m^3}$ من القشرة لكل كيلومتر من الظهرة؛ ومع كسر مصهور متوسط قدره \qty{12}{\percent}
تأتي من $\num{3.6e5}/0.12 = \qty{3.0e6}{m^3}$ من الوشاح.
\end{solution}

\begin{solution}{exo:b2:mantle-melting:11}
(أ) بعد أن تغادر كتلةٌ $t$ من A وحدةَ كتلة من السائل، يحوي السائل $0.6 -
t$ من A، و0.2 من B و0.2 من C في كتلة $1 - t$: $\big((0.6 - t), 0.2,
0.2\big)/(1 - t)$، وتبقى نسبة B إلى C مساوية 1. (ب) على \omterm{def:b2:mantle-melting:ternary}{الخط المتبلور المشترك} $(0.40 - 0.15s, 0.60 -
0.30s, 0.45s)$؛ وتساوي B وC يعطي $0.60 - 0.30s = 0.45s$، $s = 0.8$: فللسائل $L_1$ التركيب A 28، B
36، C 36. ثم $0.2/(1 - t) = 0.36$، $t = 0.44$: فقد تبلور \qty{44}{\percent} من النظام
على هيئة A. (ج) C كله في السائل: $f = 20/45 = 0.44$، أي
\qty{44}{\percent} لا يزال سائلًا. (د) A: $0.60 - 0.444\times0.25 = 0.49$؛ B: $0.20 -
0.444\times0.30 = 0.07$؛ ومع السائل، $49 + 7 + 44 = \qty{100}{\percent}$.
\end{solution}

\begin{solution}{exo:b2:mantle-melting:12}
(أ) $X_{\mathrm{sat}}(1) = 12 + 1 = \qty{13}{\percent}$، $\Delta T =
\qty{294}{\celsius}$، وخط الصلابة المشبع $1213.5 - 294 = \qty{919}{\celsius}$؛
$X_{\mathrm{sat}}(2) = 12\times1.516 + 2 = \qty{20.2}{\percent}$، $\Delta T =
43\times20.2^{0.75} = 43\times9.53 = \qty{410}{\celsius}$، وخط الصلابة المشبع $1331.1
- 410 = \qty{921}{\celsius}$. (ب) يذيب المصهور ماءً أكثر عند ضغط أعلى،
فيلغي الخفضُ الأكبر ارتفاعَ خط الصلابة الجاف بمقدار \qty{118}{\celsius}.
(ج) يجب أن ينخفض خط الصلابة بمقدار $1331.1 - 1100 = \qty{231}{\celsius}$: $X =
(231/43)^{1/0.75} = 5.37^{1.333} = \qty{9.4}{\percent}$ من الماء، أي أقل من
\qty{20.2}{\percent} عند الإشباع، فيمكن للإسفين أن ينصهر.
\end{solution}

\begin{solution}{pb:b2:mantle-melting:1}
\textbf{1.} $1350 + 273 = \qty{1623}{K}$.
\textbf{2.} $\num{4.0e-5}\times1623/(3300\times1000) = \qty{1.97e-8}{K/Pa} =
\qty{19.7}{K/GPa}$، ويُقرَّب إلى \qty{20}{K/GPa}.
\textbf{3.} $\num{e9}/(3300\times9.81) = \qty{30.9}{km}$؛ $20/30.9 =
\qty{0.65}{K/km}$.
\textbf{4.} $1350 + 20P = 1085.7 + 132.9P - 5.1P^2$، أي $5.1P^2 - 112.9P +
264.3 = 0$: $P_0 = (112.9 - \sqrt{12746.4 - 5391.7})/10.2 = (112.9 - 85.76)/10.2 =
\qty{2.66}{GPa}$.
\textbf{5.} $2.66\times30.9 = \qty{82}{km}$؛ $T = 1350 + 20\times2.66 =
\qty{1403}{\celsius}$.
\textbf{6.} $1085.7 + 132.9\times2.66 - 5.1\times2.66^2 = 1085.7 + 353.5 - 36.1 =
\qty{1403}{\celsius}$.
\textbf{7.} $T_{\mathrm s} = 1085.7 + 199.4 - 11.5 = \qty{1273.6}{\celsius}$؛
$T_{\mathrm l} = 1475.0 + 120.0 - 7.2 = \qty{1587.8}{\celsius}$؛ والفرق
\qty{314.2}{K}.
\textbf{8.} $1273.6 + 0.10\times314.2 = \qty{1305.0}{\celsius}$، أي \qty{1578}{K}.
\textbf{9.} $T_{\mathrm s}' = 132.9 - 15.3 = \qty{117.6}{K/GPa}$، $T_{\mathrm l}' =
80.0 - 9.6 = \qty{70.4}{K/GPa}$: $(\partial T/\partial P)_F = 117.6 +
0.10\times(70.4 - 117.6) = \qty{112.9}{K/GPa}$؛ $(\partial T/\partial F)_P =
\qty{314.2}{K}$.
\textbf{10.} $c_p/T = 1000/1578 = 0.634$؛ البسط $0.634\times112.9 - 12.1 = 71.6
- 12.1 = 59.5$؛ والمقام $300 + 0.634\times314.2 = 499.2$: $0.119$، أي
\qty{11.9}{\percent} لكل غيغاباسكال.
\textbf{11.} مع الدالة المضبوطة $F = T'^{1.5}$، حيث $T'$ كسر المسافة من خط الصلابة إلى خط السيولة، يجب أن ترتفع درجة الحرارة بشدّة مع $F$ لصنع
أول واحد في المئة من المصهور، فيكون $(\partial T/\partial F)_P$ كبيرًا قرب خط الصلابة
ويبدأ الانصهار ببطء. $F_{\max} = 0.09\times2.66 = 0.24$، أي \qty{24}{\percent}.
\textbf{12.} $0.13/0.50 = 0.26$: يُستنفد الكلينوبيروكسين عند \qty{26}{\percent} من
الانصهار. وعند \qty{24}{\percent} لا تزال البقايا تحت المحور تحوي قليلًا منه:
فهي لرزوليت فقير بالكلينوبيروكسين، على حافة \omterm{def:b2:mantle-melting:peridotites}{الهارزبورجيت}.
\textbf{13.} $h_c = 0.24\times82/2 = \qty{9.8}{km}$.
\textbf{14.} $F_{\max}/2 = \qty{12}{\percent}$.
\textbf{15.} \qty{9.8}{km} أكثر بنسبة \qty{38}{\percent} من \qty{7.1}{km}، أي فوق
مدى القشرة العادية.
\textbf{16.} الإنتاجية الثابتة تعطي الوشاح الأعمق مصهورًا أكثر من اللازم؛ أمّا
النموذج المضبوط فيصنع المصهور الأول ببطء، فتبقى $F(z)$ أدنى قرب البدء،
ويكون $\int F\,\mathrm dz$ أصغر.
\textbf{17.} $F_{\max} = \Pi P_0$ و$H = 30.9P_0$ (بالكيلومتر): $h_c = \Pi P_0\times30.9P_0/2
= kP_0^2$ مع $k = 0.09\times30.9/2 = \qty{1.39}{km/GPa^2}$.
\textbf{18.} $P_0 = \sqrt{7.1/1.39} = \sqrt{5.11} = \qty{2.26}{GPa}$.
\textbf{19.} $2.26\times30.9 = \qty{70}{km}$.
\textbf{20.} $T_{\mathrm s}(2.26) = 1085.7 + 300.4 - 26.0 = \qty{1360.0}{\celsius}$،
و$T_p = 1360.0 - 20\times2.26 = 1360.0 - 45.2 = \qty{1315}{\celsius}$.
\textbf{21.} من أجل \qty{6.3}{km}: $P_0 = \sqrt{4.53} = \qty{2.13}{GPa}$، $T_{\mathrm s}
= \qty{1345.6}{\celsius}$، $T_p = \qty{1303}{\celsius}$. ومن أجل \qty{7.9}{km}: $P_0 =
\sqrt{5.68} = \qty{2.38}{GPa}$، $T_{\mathrm s} = \qty{1373.1}{\celsius}$، $T_p =
\qty{1326}{\celsius}$. فتشتّت القشرة يسمح بالقيمة $1315 \pm
\qty{12}{\celsius}$.
\textbf{22.} $P_0 = \sqrt{20/1.39} = \sqrt{14.39} = \qty{3.79}{GPa}$؛ $T_{\mathrm s} =
1085.7 + 503.7 - 73.3 = \qty{1516.1}{\celsius}$؛ $T_p = 1516.1 - 75.8 =
\qty{1440}{\celsius}$، أي زيادة قدرها $1440 - 1315 = \qty{125}{\celsius}$.
\textbf{23.} يستعمل الطريقان شواهد مختلفة. فالحمم تسجّل أسخن وشاح،
في مركز العمود؛ أمّا القشرة فتأخذ متوسط مصاهير كل الوشاح الذي
يعبر منطقة الانصهار، وكثير منه أبرد. والنموذج البسيط يُبقي أيضًا
الإنتاجية ثابتة، مع أن الكلينوبيروكسين يُستنفد عند \qty{26}{\percent}
بينما $\Pi P_0 = 0.09\times3.79 = 0.34$ هنا: فالوشاح الساخن الحقيقي ينصهر أقل، ويحتاج إلى
درجة حرارة أعلى ليصنع القشرة نفسها.
\textbf{24.} $T_{\mathrm s}'(2.26) = 132.9 - 23.1 = \qty{109.8}{K/GPa}$؛ $\Delta P_0 =
10/(109.8 - 20) = \qty{0.111}{GPa}$؛ وبما أن $h_c \propto P_0^2$، فإن $\Delta h_c/h_c =
2\times0.111/2.26 = \qty{9.8}{\percent}$، أي نحو العُشر لكل \qty{10}{\celsius}.
\textbf{25.} عند ظهرة بطيئة جدًا يصعد الوشاح ببطء شديد حتى يفقد حرارته
بالتوصيل نحو السطح البارد؛ فيغادر التدرّج الكظوم، ويتوقف انصهاره
أعمق، ويُقتطع رأس المثلث، حيث $F$ أكبر ما يكون: فتكون القشرة
أرقّ.
\textbf{26.} تعطي قشرة الظهرات العادية درجة حرارة كامنة للوشاح قدرها
\textbf{\qty[text-series-to-math=true]{1315}{\celsius}}، في حدود نحو
\qty{12}{\celsius}: أي أدنى بمقدار \qty{35}{\celsius} من \qty{1350}{\celsius} في فيزياء
المعادن، وهي غير مؤكدة بالقدر نفسه تقريبًا، وداخل المدى من 1280 إلى
\qty{1400}{\celsius} من الحمم بوضوح. فثلاثة مقاييس حرارة مستقلة تتفق على وشاح
قرب \qty{1300}{\celsius} إلى \qty{1350}{\celsius}.
\end{solution}
